\documentclass[a4paper,12pt]{article}
\usepackage{fontspec}
\usepackage{unicode-math}
\setmathfont{Latin Modern Math}
\usepackage{amsmath,amssymb}
\usepackage{graphicx}
\usepackage[hidelinks]{hyperref}
\usepackage[margin=2.5cm]{geometry}
\usepackage[numbers]{natbib}
\usepackage{booktabs,array,float}
\title{From Compute to Complexity}
\author{Qinrang Liu}
\date{June 2026}
\begin{document}
\maketitle
\begin{abstract}
The rapid scaling of large language models has delivered remarkable functional capabilities yet produced exponentially growing energy costs with sub-linear returns---a thermodynamic trajectory that converges not toward general intelligence but toward an unsustainable asymptote. We argue that this trajectory is not an engineering deficiency but a consequence of pursuing the wrong variable: compute, rather than complexity. Von Neumann identified in 1948 that intelligence requires a complexity threshold; here we quantify that threshold through a framework grounded in thermodynamic phase transitions, renormalization group theory, and complex network science. The result is the Coordination Spatiotemporal Complexity theorem: CST = (Sc $\cdot$ Tc) $\cdot$ exp($\alpha$ $\cdot$ $\Gamma$st), where structural integration, dynamical richness, and their physical coupling jointly determine emergent intelligence potential. We derive six universal thresholds at natural constants {1/$\sqrt{2}$, 1, $\phi$, e, $\pi$, $\delta$} and validate across 40 biological and artificial systems spanning 8 taxonomic grades and 20 distinct ANN/NMH architectures (Spearman $\rho$ = 0.976, 100\% accuracy under UCCP normalization). Neuromorphic hardware (Intel Loihi-2) is separately classified from binary-digital ANN, confirming the $\alpha$-barrier prediction. Intelligence Efficiency $\eta$\_I reveals an approximately six-order-of-magnitude gap between brains and current AI, and a four-generation hardware roadmap identifies the physically necessary path from present systems to general intelligence. (150 words) Keywords: intelligence emergence; complexity threshold; von Neumann; spatiotemporal coordination; intelligence efficiency; phase transitions; neuromorphic computing
\end{abstract}

\section{Introduction}

The sustainability crisis of artificial intelligence. The trajectory of modern AI development is defined by a single operating principle: scale compute, and intelligence will follow. Each generation of frontier LLMs has required substantially greater training compute than its predecessor, with scaling law analyses projecting continued exponential growth [31]. Inference energy has grown proportionally. Yet empirical scaling laws now reveal that capability improvements per unit energy expenditure follow a sub-linear curve---each successive generation buys less intelligence per joule invested. The global AI industry is approaching a thermodynamic asymptote---one enforced not by CMOS fabrication technology per se, but by the binary digital logic paradigm implemented on it: the current paradigm can produce ever more capable functional systems, but the energy cost required to sustain them grows without bound while the gap between these systems and genuine general intelligence does not close.


This is not merely a resource problem. It is a symptom of pursuing the wrong quantity. The dominant paradigm equates intelligence with compute---more parameters, more data, more hardware---and measures progress by benchmark performance. But benchmark performance and intelligence emergence are orthogonal dimensions. GPT-class models surpass most humans on standardized tests in law, medicine, and coding. Yet as we show below, GPT-2---a representative large-scale open-weight language model---scores approximately 30-fold lower than the human brain on the metric of emergent intelligence potential (CST = 0.056 vs. 3.909), and even below Caenorhabditis elegans, a 279-neuron nematode (CST = 0.357 under correct graded-potential physics). This is not a contradiction. It is a revelation: we have been measuring the wrong thing.


The von Neumann threshold and the complexity imperative. The foundations for a different view were laid before modern AI existed. Von Neumann, in his 1948 lectures on the theory of self-reproducing automata [44] (published 1966)---building on the computational foundations laid by Turing [45]--- identified a critical complexity threshold below which systems can only simplify and above which genuine self-organization and reproduction become possible. This threshold was not defined by computational power but by structural and dynamical complexity---the richness of a system's internal organization. The insight was prophetic but remained qualitative for seven decades: how to measure this complexity, and what its quantitative thresholds are, were open questions.


The intervening decades produced fragments of an answer. Criticality theory showed that neural systems operate near phase transitions [6,7], where small changes in network state produce disproportionate changes in dynamics---a signature of complexity at the edge of chaos [50]. This dynamical framework has since been formalized by the phenomenological renormalization group [51], revealing that scale-invariant criticality in neural tissue is not an approximation but a universal phase, with each coarse-graining step preserving the statistical structure of neural correlations---directly underpinning the exponential coupling term in CST (see Theory). Complex network theory revealed that biological neural networks share universal structural properties: small-world topology [8], hierarchical modularity [9], and broad degree distributions with hierarchical organization [48,49]---properties that distinguish them from the uniform-connectivity graphs of artificial neural networks. Thermodynamic analysis of information processing showed that physical coupling between structure and function---not just the existence of structure or function separately---is what distinguishes adaptive from reflexive behavior [23]. Intelligence itself has been argued to be intrinsically dynamical rather than representational: emergent coherent order arising from local nonlinear interactions under physical constraints [52], a characterization that directly maps onto the CST formalism.


From fragments to a unified theory. The present work assembles these fragments into a single quantitative framework by asking: what is the minimal set of physical quantities whose joint optimization is both necessary and sufficient for intelligence emergence? The answer, derived from first principles rather than fitted to data, is three quantities and their interaction: spatial network complexity Sc (how richly connected and hierarchically organized a network is), temporal dynamical complexity Tc (how rich and multi-timescale the network's spontaneous dynamics are), and crucially, the coupling $\Gamma$st between them---the degree to which the network's functional dynamics are physically aligned with its structural organization.


The critical insight is that these quantities do not add; they multiply and amplify. A network with rich structure and poor dynamics, or rich dynamics and poor structure, achieves modest complexity. But when structure and function are physically coupled, each reinforces the other in a cascade process formally equivalent to information gain near a phase transition [6]. This is why the coupling term enters the equation exponentially: CST = (Sc $\cdot$ Tc) $\cdot$ exp($\alpha$ $\cdot$ $\Gamma$st). The coefficient $\alpha$ = ln($M_{\text{eff}}$) is determined entirely by device physics---the number of distinguishable states a node can occupy---making it the one variable that hardware, not software, controls absolutely.


The six intelligence thresholds {1/$\sqrt{2}$, 1, $\phi$, e, $\pi$, $\delta$} are not empirically fitted; they are derived from the symmetry-breaking structure of phase transitions in complex networks, in the same mathematical tradition that gives thermodynamics its universal constants. Their validation across 40 biological and artificial systems---with no free parameters---is the empirical test of a physical theory, not a data-fit.


Existing frameworks address fragments of this picture [1--3]: Integrated Information Theory (IIT) proposes $\Phi$ as a consciousness measure [4], but computation scales as O(2\textsuperscript{n}), limiting it to \textasciitilde{}30 nodes [5]; criticality theory does not predict intelligence levels [6,7]; complex network theory lacks a unified metric connecting structure to emergent behavior [2,9]. The CST framework provides the unification.


We further show that the global AI industry's architectural evolution over 2017--2025 constitutes independent empirical validation: every major architectural innovation---from MoE modularity and NAS-optimized hierarchy, to SSM recurrence and continuous-time liquid dynamics, to inference-time plasticity---maps onto a specific CST component, confirming that the industry has empirically converged toward CST-optimal architecture through engineering pressure alone, while simultaneously revealing the one transition the scaling paradigm cannot make: from simulated $\Gamma$st to physical $\Gamma$st.



\section{Results}

We formalize the CST theorem on five axioms. These are not arbitrary postulates.
From these axioms:
\begin{equation}
CST = (S_c \cdot T_c) \cdot \exp(\alpha \cdot \Gamma_{st}) \tag{1}
\end{equation}
Spatial complexity Sc quantifies structural integration potential.


\end{document}